% Volume VI: Adversaries, Platforms, and Automated Judgment
% Part XII. Platforms and Automated Publicity
% Chapter 70: Synthetic Evidence and Provenance Failure
% Spine: proof-spines/ch070-spine.md

\chapter{Synthetic Evidence and Provenance Failure}
\label{ch:ch070}


Generative media raises the cost of verification, but it does not make all evidence worthless. The problem is better described as \textbf{provenance failure}: an environment in which artifacts can be fabricated, remixed, stripped of origin, and recirculated faster than ordinary trust practices can stabilize their history.

The chapter's technical inventory is consequently layered rather than magical. Cryptographic provenance, sensor authentication, chain of custody, watermarking, and model-based detection can each strengthen a record's status, while each also has limits that prevent any purely technical guarantee. Synthetic evidence therefore shifts emphasis from naive faith in appearances to auditable pathways by which an image, clip, or document acquires trustworthy lineage. The task is not to declare all media suspect, but to rebuild evidentiary confidence around provenance systems that remain contestable and incomplete. When institutions demand immaculate provenance even where authentic witnesses cannot supply it, evidentiary rigor flips into an unmeetable burden that excludes real evidence along with synthetic fraud.
